\documentclass[a4paper,11pt]{article}

% Deze template van Johan Vervloet en Wim Willems
% is gelicenseerd onder een Creative
% Commons Naamsvermelding-GelijkDelen 3.0 Unported
% licentie (http://creativecommons.org/licenses/by-sa/3.0/)

% Nodige packages voor compilatie:
%  Fedora:
%    texlive-textpos, texlive-multirow, texlive-pdftex-def, texlive-SIunits
%  Debian:
%    texlive-latex-extra

\usepackage[dutch]{babel}
\usepackage[pdftex]{graphicx}
\usepackage{multirow}
\usepackage[absolute]{textpos}

\usepackage{hyperref}                   % Om hrefs te maken
\usepackage{url}                        % Om url's te verwerken
\usepackage{graphicx}                   % Om figuren te kunnen verwerken
\usepackage{xspace}                     % Magische spaties na een commando
\usepackage[utf8]{inputenc}             % Om niet ascii karakters rechtstreeks te kunnen typen
\usepackage{float}                      % Om nieuwe float environments aan te maken. Ook optie H!
\usepackage{flafter}                    % Opdat floats niet zouden voorsteken
\usepackage{listings}                   % Voor het weergeven van letterlijke text en codelistings
\usepackage{eurosym}                    % om het euro symbool te krijgen
\usepackage{textcomp}                   % Voor onder andere graden celsius
\usepackage[Gray,squaren,thinqspace,thinspace]{SIunits} % Om elegant eenheden te zetten


\setlength{\TPHorizModule}{25mm}
\setlength{\TPVertModule}{\TPHorizModule}
\textblockorigin{15mm}{15mm}

% Sans-serif font, volgens huisstijl
%
\renewcommand{\familydefault}{\sfdefault}

\begin{document}

%% algemene instellingen voor listings

%%
%% Eerst een blok rommel voor de 'huisstijlheader'
%%

\begin{textblock}{2}(0,0)
\includegraphics[height=32mm]{chiro.png} \\
\end{textblock}

\begin{textblock}{4.1}(3,0)
\textbf
{
\begin{tabular}{|p{65mm}|p{10mm}|r|}
\multicolumn{3}{r}{\Large Chirojeugd-Vlaanderen vzw} \\
\multicolumn{3}{r}{\,} \\
\hline
\multicolumn{2}{|l|}{Werkgroep Informatica} & 271 \\
\hline
25 april 2008 & Doc. & 2 \\
\hline
\end{tabular} 
}
\end{textblock}

\vspace*{4mm}

%%
%% HIER BEGINT HET DOCUMENT
%%


\section{Eerste sectie}

Blabli blabla. Tralali tralala. Jaja.

\end{document}
